\section[Rédaction Bourbaki \redac{398} : dualité des modules de longueur finie]{Rédaction Bourbaki \redac{398} : dualité des modules de longueur finie sur un anneau principal\protect\verifie{Redaction398}}\label{sec:bourbaki-398}

La rédaction \redac{398} des archives Bourbaki (Archives Henri Poincaré, \ahp{949}), datée d'avril 1963, est une contre-rédaction : elle reprend le sujet de la rédaction \redac{397} (section~\ref{sec:bourbaki-397}), compléments au chapitre VII d'\emph{Algèbre}, mais en s'interdisant les modules injectifs et en se limitant aux modules de longueur finie. Soient $A$ un anneau principal, $K$ son corps des fractions, supposé distinct de $A$, et, pour un $A$-module $E$, $E^{*} = \operatorname{Hom}_A(E, K/A)$. L'évaluation $(x, x^{*}) \mapsto x^{*}(x)$ est bilinéaire et fournit une application canonique $i_E \colon E \to E^{**}$. Le texte énonce trois propositions : l'autodualité des modules de longueur finie (proposition~1), l'exactitude du foncteur $E \mapsto E^{*}$ (proposition~2) et la bijection entre sous-modules de $E$ et de $E^{*}$ par passage à l'orthogonal (proposition~3) ; il signale enfin le cas $A = \mathbf Z$, $K/A = \mathbf Q/\mathbf Z$, qui est la dualité des groupes abéliens finis. Une main non identifiée a ajouté en marge la base duale explicite. Sa transcription reste locale pour l'instant ; on n'en cite ici aucun passage.

\emph{L'idée.} Le théorème de structure ramène tout au module cyclique $A/tA$, $t \neq 0$. Une forme linéaire sur $A/tA$ est déterminée par l'image $q$ de la classe de $1$, sous la seule condition $tq = 0$ ; le dual de $A/tA$ est donc le sous-module $t^{-1}A/A$ de $K/A$, et $a \mapsto a/t \bmod A$ l'identifie à $A/tA$.

\begin{encadre}
\begin{enonce}[Bourbaki, rédaction \redac{398}, propositions~1 à 3]
Soient $A$ un anneau principal, $K$ son corps des fractions et $E$ un $A$-module de longueur finie (c'est-à-dire de type fini et de torsion).
\begin{enumerate}
\item[a)] Pour $t \neq 0$, $(A/tA)^{*}$ s'identifie à $t^{-1}A/A$, donc est isomorphe à $A/tA$.
\item[b)] $E^{*}$ est isomorphe à $E$ (non canoniquement en général), et $i_E \colon E \to E^{**}$ est un isomorphisme.
\item[c)] Si $E' \to E \to E''$ est une suite exacte, la suite $E''^{*} \to E^{*} \to E'^{*}$ est exacte.
\item[d)] L'application $F \mapsto F^{0}$, où $F^{0} \subset E^{*}$ est l'orthogonal de $F$, est une bijection décroissante de l'ensemble des sous-modules de $E$ sur celui des sous-modules de $E^{*}$ ; elle échange sommes et intersections.
\end{enumerate}
\end{enonce}
\end{encadre}

\begin{proof}
a) L'évaluation en la classe de $1$ est injective sur $(A/tA)^{*}$, et son image est formée des $q \in K/A$ tels que $tq = 0$ : si $tq = 0$, l'application $a \mapsto aq$ s'annule sur $tA$ et passe au quotient. Ce sous-module est aussi l'image de $\delta_t \colon a \mapsto a/t \bmod A$ : si $q$ est la classe de $k$ et $tk = b \in A$, alors $q = \delta_t(b)$. Enfin $\delta_t(a) = 0$ si et seulement si $a/t \in A$, c'est-à-dire $a \in tA$ ; donc $(A/tA)^{*} \simeq A/tA$. Explicitement, et c'est la formule de la note marginale, la forme qui correspond à la classe de $1$ envoie la classe de $1$ sur la classe de $t^{-1}$ dans $K/A$ ; elle engendre $(A/tA)^{*}$. Pour $E = \bigoplus_i A/t_iA$, somme finie avec $t_i \neq 0$, de base les classes $e_i$ de $1$, les formes $e_i^{*}$ définies par $e_i^{*}(e_j) = \delta_{ij}\, t_i^{-1} \bmod A$ forment la base duale : elles sont l'image des $e_i$ par un isomorphisme $E \simeq E^{*}$.

b) Par le théorème de structure, $E \simeq \bigoplus_i A/p_i^{e_i}A$, somme finie. Le dual d'une somme directe finie est le produit des duaux, et a) s'applique à chaque facteur ; d'où $E^{*} \simeq E$. Pour $i_E$ : $K/A$ est divisible, donc injectif (critère de Baer, voir la rédaction \redac{397}). Si $x \neq 0$ a pour annulateur $(s)$, l'élément $s$ n'est pas inversible, et l'application $ax \mapsto a/s \bmod A$ de $Ax$ dans $K/A$ est bien définie et ne s'annule pas en $x$ ; elle se prolonge à $E$, d'où l'injectivité de $i_E$. Comme $E^{**} \simeq E^{*} \simeq E$, la composée de $i_E$ et d'un tel isomorphisme est un endomorphisme injectif de $E$ ; $E$ étant artinien, il est surjectif, donc $i_E$ l'est aussi.

c) L'inclusion de l'image dans le noyau est immédiate. Si $\varphi \in E^{*}$ s'annule sur l'image de $E'$, c'est-à-dire sur le noyau de $g \colon E \to E''$, elle se factorise par $E/\ker g \simeq g(E)$, et la forme obtenue sur $g(E)$ se prolonge à $E''$ par injectivité de $K/A$.

d) Pour $x \notin F$, l'injectivité de $i_{E/F}$ fournit une forme sur $E/F$ non nulle en la classe de $x$, donc un élément de $F^{0}$ non nul en $x$ : ainsi l'orthogonal de $F^{0}$ dans $E$ est $F$. Pour $G \subset E^{*}$ et $\varphi \notin G$, le même argument appliqué à $E^{*}$ (de longueur finie par b)) donne $\Phi \in E^{**}$ nulle sur $G$ et non nulle en $\varphi$ ; comme $i_E$ est surjectif, $\Phi = i_E(x)$, et $x$ est orthogonal à $G$ sans l'être à $\varphi$. Les deux passages à l'orthogonal sont donc des bijections réciproques et décroissantes. L'orthogonal d'une somme est toujours l'intersection des orthogonaux ; en appliquant ceci dans $E^{*}$ et la bijectivité, l'orthogonal d'une intersection est la somme des orthogonaux.
\end{proof}

\noindent La contre-rédaction évite à dessein les modules injectifs : elle obtient b) par réduction au cas cyclique et c) par un décompte des longueurs. La preuve ci-dessus, comme la formalisation, utilise au contraire l'injectivité de $K/A$, plus économique en Lean ; elle montre au passage que c) vaut sans hypothèse de finitude et que $i_E$ est injectif pour tout module de torsion.

\begin{sloppypar}
\noindent\emph{Vérification en Lean.} Le module $K/A$ est \texttt{KModA} et le dual $E^{*}$ est \texttt{Dual}. a) est \texttt{range\_evalOne} et \texttt{dualQuotEquiv} (avec \texttt{ker\_divBy} et \texttt{range\_divBy}), la formule explicite \texttt{dualQuotEquiv\_symm\_mk\_one\_apply} (avec \texttt{dualGen} et \texttt{dualGen\_mk\_one}), et la base duale de $\bigoplus_i A/t_iA$ \texttt{dualBasis}, avec \texttt{dualBasis\_lof}, \texttt{dualSumEquiv} et \texttt{dualSumEquiv\_symm\_lof} ; b) \texttt{nonempty\_dual\_equiv}, puis \texttt{evalDual\_injective}, \texttt{evalDual\_bijective} et l'isomorphisme \texttt{evalDualEquiv} ; c) \texttt{dual\_exact} et \texttt{dual\_surjective\_of\_injective}, à partir de \texttt{baer} ; d) \texttt{orthOrderIso}, \texttt{orth\_sup} et \texttt{orth\_inf}, dans \href{\leanblobmain/Redaction398.lean}{\texttt{Redaction398.\allowbreak{}lean}}. Un module de longueur finie y est codé comme un module de type fini et de torsion, et le théorème de structure utilisé est celui de mathlib (\texttt{Module.equiv\_directSum\_of\_isTorsion}).
\end{sloppypar}

\litterature{retrouvée. C'est le cas d'un anneau principal de la dualité de Matlis (Matlis, \emph{Pacific J. Math.} 8, 1958 ; Bruns--Herzog, \emph{Cohen--Macaulay rings}, §~3.2), $K/A$ étant l'enveloppe injective de la somme des $A/\mathfrak m$ ; pour $A = \mathbf Z$, c'est la dualité classique des groupes abéliens finis. mathlib n'avait que ce dernier cas, à valeurs dans $\mathbf C$ (\texttt{AddChar.doubleDualEquiv}), et non l'énoncé sur un anneau principal quelconque ; c'est ce qui justifie la formalisation. L'édition ne le donne pas pour nouveau.}
